\section{Alternative substrates and the limits of the present search}
\label{sec:alternative-substrates}

The construction problem does not logically require a particular Turing
machine, nor a particular universal-polynomial presentation.  We examined
alternatives to the inherited substrate as well as algebraic changes to its
final arithmetic kernel.  None of the independent-substrate routes examined
here produced a smaller, fully instantiated polynomial with a complete
correctness argument.  This is a statement about the work completed in this
iteration, not an impossibility theorem or a claim that the existing
substrate is optimal.

\subsection{A direct proof-checker route}

Fix an effective proof calculus for ZFC and an explicit coding of its
formulas and finite proofs.  The predicate
\(\operatorname{Ref}_{\mathrm{ZFC}}(p)\), asserting that \(p\) codes a
refutation, can be chosen primitive recursive.  A direct compiler could
therefore replace the machine simulation by a polynomial
\(D(p,\mathbf x)\) with a verified representation theorem
\[
 \mathrm{ZFC}\vdash
 \bigl(\exists\mathbf x\in\mathbb N^m\;
       D(p,\mathbf x)=0\bigr)
 \longleftrightarrow \operatorname{Ref}_{\mathrm{ZFC}}(p).
\]
After existentially quantifying \(p\), the absence of zeros would be
equivalent to syntactic \(\operatorname{Con}(\mathrm{ZFC})\).  This is a
valid architecture that makes no reference to K215.

To turn that architecture into the requested fixed equation, the following
data and arguments must actually be supplied.  The syntax encoding must
specify formula parsing, bound-variable handling and substitution; axiom
recognition must include the Separation and Replacement schemes; and the
proof checker must specify every allowed inference.  The arithmetic
compiler must then emit polynomial syntax, literal fixed coefficients and
a proof relating its integer witnesses to the checker.  Finally, the
checker must be calibrated against the same effective ZFC calculus used in
the consistency assertion.  It is the combined cost of these components
that should be compared with the present construction.

The MRDP theorem establishes the availability of a Diophantine
representation, but an unspecified representation or an unspecified
universal-polynomial parameter is not an explicit coefficient instance.
Likewise, a theorem accepting supplied representation syntax does not by
itself construct that syntax for a particular proof checker.  The current
ProveIt \texttt{QUARTIC\_COMPILER.md} explicitly distinguishes these two
interfaces.  Closely related distinctions between mathematical
construction, polynomial data and formal verification also occur in the
Isabelle universal-pair development.\footnote{Jonas Bayer and Marco David,
\emph{A Formal Proof of Complexity Bounds on Diophantine Equations},
\url{https://arxiv.org/abs/2505.16963}; the associated development is
\url{https://isa-afp.org/entries/Diophantine_Universal_Pairs.html}.}

A finite logical presentation or a smaller equiconsistent theory could be
a useful source language.  However, counting axioms does not measure the
size of the compiled polynomial.  Such a proposal must provide effective
proof translations sufficient to justify the intended consistency
equivalence, and must still pay for its syntax checker.  Replacing proofs
by finite models of ZFC would not work: ZFC itself proves that its Infinity
axiom rules out finite models.  Replacing the theory by an assertion that
a model of ZFC exists also requires a separate strength analysis; that
assertion cannot simply be treated as an inexpensive equivalent axiom.

\subsection{Generic degree reduction is not a new small quartic}

An arithmetic circuit can be converted to quadratic equations by
introducing values for intermediate registers, and a sum of their squares
then has degree at most four.  This exchanges degree for unknowns.  Our
straightforward lift of the ordinary Pell kernel did not improve the
existing source-31 point: that source already has 31 positive unknowns and
22 quadratic residuals.  Its comparison evaluator has 60 operations and
its quartic sum of squares has 125 operations.  These figures are the
existing \texttt{31-sun10/quartic31/QUARTIC60\_PROOF.md} benchmark, not a
new result of this report.

One illustrative failed elimination concerns \(X=wq\).  Combining it
with the transport relation can remove the letter \(w\) from a polynomial
identity, but generally leaves only \(q\mid XC\).  Restoring the old
witness \(w=X/q\) would require \(q\mid X\).  It is invalid to cancel
\(C\) without first establishing \(\gcd(q,C)=1\).  In the contemplated
route, the desired coprimality depended on source-decoding consequences
that themselves used \(q\mid X\).  We therefore did not accept this as a
reduction in the number of variables.  This identifies a missing argument
in that particular route; it does not rule out a different elimination.

\subsection{Two auxiliary congruence routes that do not lower degree}

The degree-24 integral-root construction has an auxiliary Pell norm
\[
 U V^2-(U-1)y^2=1,
 \qquad U=I^2=\Delta(f^2-1),
 \qquad y=V+u.
\]
Here \(\deg U=6\), \(\deg V=5\), and the supplied gap \(u\) has degree
one.  The gap residual
\[
 (V+u)^2-Uu(2V+u)-1
\]
has degree 12, giving degree 24 after squaring.  We examined whether
canonical congruences permit a cheaper replacement of the gap.

For the canonical odd Pell index \(R\equiv3\pmod4\),
\(y=\psi_I(R)\equiv-1\pmod{I^2}\).  Thus
\(t=(y+1)/U\) is a positive integer.  The exact identity
\begin{align*}
 &U V^2-(U-1)(Ut-1)^2-1\\
 &\hspace{12mm}=
 U\bigl[V^2-U(U-1)t^2+2(U-1)t-1\bigr]
\end{align*}
gives a valid alternative residual when completeness uses these canonical
witnesses.  Its degree is 14, because of \(U(U-1)t^2\), so its square has
degree 28.  This useful divisibility does not lower the degree of the
current gap construction.  Supplying \(Uu\) instead leads to denominator
clearing and likewise did not improve the measured degree.

There is a second canonical congruence.  Write
\(A=a+2\), so \(\Delta=A^2-1\), and let
\(D=\chi_A(R)\), \(c=\psi_A(R)\).  The odd-subsequence Pell
identities, together with
\(I^2\equiv1-A^2\pmod{f^2}\), give
\[
 Ay\equiv-D\pmod{f^2}
 \qquad(R\equiv3\pmod4).
\]
For example, at \(R=3\) one has
\(\psi_I(3)=4I^2-1\), and multiplication of its reduction modulo
\(f^2\) by \(A\) gives \(-\chi_A(3)\).  The general identity follows
from the same odd-subsequence recurrence.  Notice that reducing the norm
alone gives only \((Ay)^2\equiv D^2\pmod{f^2}\); the negative branch
cannot be selected for arbitrary supplied solutions from that square
congruence alone.

The canonical congruence suggests
\(y=(f^2t-D)/A\).  Write \(V=f^2W-c\) and center the new quotient by
\(t=AW+v\).  Formally retaining \(D\) as a register, after clearing the
denominator \(A^2\) and using \(D^2-\Delta c^2=1\), the norm numerator
is divisible by \(f^2\).  The resulting expression still contains
\(-2A\Delta Wv f^4\).  This term has degree 14 under
\(\deg A=2\), \(\deg\Delta=4\), \(\deg W=3\) and
\(\deg v=\deg f=1\).  Moreover, source 58 reconstructs \(D\) rather
than supplying it as a separate unknown, so that reconstruction would also
have to be accounted for in a literal polynomial.  This calculation does
not provide the hoped-for degree-22 equation.  It proves no broader
obstruction to using this congruence in a different construction.

\subsection{Scope of the reported advance}

The demonstrated improvements in this report concern the explicit
constructions and comparisons proved in its preceding sections.  The
alternative proof substrates above remain research proposals.  They are
not counted as delivered equations, and the calculations that failed to
improve degree are retained to make the search boundary clear.  In
particular, we make neither a global simplicity-minimum claim nor a claim
that a construction independent of K215 or Jones is impossible.

Throughout, a proof of
\(\mathrm{ZFC}\vdash[\text{no zeros}\leftrightarrow
\operatorname{Con}(\mathrm{ZFC})]\)
is an internal equivalence theorem.  It is not a proof in ZFC of either
side.  Finite coefficient, identity and degree audits establish their
stated finite properties; they do not establish the consistency assertion
by failing to find a zero.
